\chapter{Arithmetic and Geometric Sequences and Series}
\label{ch:ch18-sequences-series}

\section{From Change to Accumulation}

\begin{historicalnote}{The Word ``Calculus''}
The word \emph{calculus} is Latin for a small pebble or counting stone.
Before written numerals were widespread, quantities were often tracked
by correspondence rather than by symbols: a shepherd might drop one
pebble into a pouch for each sheep passing through a gate, later
reversing the process to confirm that every sheep had returned, with the
pebbles serving as an external record of a quantity no numeral was
needed to name. Arranging such pebbles into grooves or columns, as on a
Roman counting board or a later abacus, let position carry meaning as
well, an early form of the place-value systems that Chapter~0 discussed
in connection with numeral systems more generally. The word survived
this entire history essentially unchanged: by the time Newton and
Leibniz developed systematic methods for computing rates of change and
accumulated totals, ``calculus'' already meant, broadly, any organized
method of calculation, which is why their new subject was named the
differential and integral calculus rather than something entirely new.
The sequences and series developed in this chapter, built from
accumulating one term after another, are in a real sense the direct
descendants of the shepherd's pebbles: both are methods for turning many
small contributions into a single total.
\end{historicalnote}

Chapter~\ref{ch:ch17-calculus-connections} studied derivatives and the
mathematics of infinitely small change: a derivative is obtained by
examining what happens as a quantity approaches zero. This chapter
introduces a different encounter with infinity. Instead of asking what
happens as a change becomes smaller and smaller, it asks what happens as
more and more terms are added together. These two ideas --- infinitely
small change and infinitely many accumulated terms --- are among the
central themes of calculus, and this chapter develops the second in a
setting that requires only algebra and the geometric ideas already
established throughout this book.

\section{Sequences}

\begin{definition}
A \emph{sequence} is a function whose domain is the positive integers.
Its terms are commonly denoted $a_1, a_2, a_3, \dots$, where $a_n$ is
called the $n$th term.
\end{definition}

A sequence may be defined explicitly or recursively.

\begin{example}
The sequence $a_n = 2n+1$ has terms $3, 5, 7, 9, \dots$ and is defined
explicitly.
\end{example}

\begin{example}
The sequence defined by $a_1 = 1$, $a_{n+1} = a_n+2$ produces the same
terms $1, 3, 5, 7, \dots$ but is defined recursively.
\end{example}

\section{Arithmetic Sequences and Series}

An arithmetic sequence is formed by repeatedly adding a fixed quantity.

\begin{definition}
An \emph{arithmetic sequence} has a constant difference $d$ between
consecutive terms: $a_{n+1} = a_n+d$, called the \emph{common
difference}.
\end{definition}

\begin{theorem}
The $n$th term of an arithmetic sequence is $a_n = a_1+(n-1)d$.
\end{theorem}

\begin{proof}
Starting from $a_1$, each new term adds one additional copy of $d$; after
$n-1$ such additions, $a_n = a_1+(n-1)d$.
\end{proof}

A sum of terms from an arithmetic sequence is called an arithmetic
series.

\begin{theorem}
The sum of the first $n$ terms of an arithmetic sequence is $S_n =
\dfrac{n}{2}(a_1+a_n)$.
\end{theorem}

\begin{proof}
Write the sum forward and backward, $S_n = a_1+a_2+\cdots+a_n$ and $S_n =
a_n+a_{n-1}+\cdots+a_1$. Adding corresponding terms gives $2S_n =
n(a_1+a_n)$, and dividing by $2$ yields $S_n = \tfrac{n}{2}(a_1+a_n)$.
\end{proof}

\begin{example}
Find the sum $1+3+5+\cdots+99$. Since $a_1=1$, $a_n=99$, and $n=50$,
$S_{50} = \tfrac{50}{2}(1+99) = 2500$.
\end{example}

\section{Geometric Sequences}

An arithmetic sequence grows by repeated addition; a geometric sequence
grows by repeated multiplication.

\begin{definition}
A \emph{geometric sequence} has a constant ratio $r$ between consecutive
terms: $a_{n+1} = ra_n$, called the \emph{common ratio}.
\end{definition}

\begin{theorem}
The $n$th term of a geometric sequence is $a_n = a_1r^{n-1}$.
\end{theorem}

\begin{proof}
Each successive term contributes one additional factor of $r$: $a_2 =
a_1r$, $a_3 = a_1r^2$, $a_4 = a_1r^3$, and continuing this pattern gives
$a_n = a_1r^{n-1}$.
\end{proof}

\begin{example}
The sequence $2, 6, 18, 54, \dots$ has common ratio $r=3$, with explicit
formula $a_n = 2\cdot3^{n-1}$.
\end{example}

\section{The Finite Geometric Series}

The most important result of this chapter is the formula for a finite
geometric sum.

\begin{theorem}
For $r \neq 1$, $a+ar+ar^2+\cdots+ar^{n-1} = \dfrac{a(1-r^n)}{1-r}$.
\end{theorem}

\begin{proof}
Let $S_n = a+ar+ar^2+\cdots+ar^{n-1}$. Multiplying by $r$ gives $rS_n =
ar+ar^2+\cdots+ar^n$. Subtracting, $S_n - rS_n = a - ar^n$, so $(1-r)S_n
= a(1-r^n)$, and dividing by $1-r$ yields $S_n = \dfrac{a(1-r^n)}{1-r}$.
\end{proof}

\begin{example}
Compute $1+2+4+8+16$. Here $a=1$, $r=2$, $n=5$, so $S_5 =
\dfrac{1-2^5}{1-2} = 31$.
\end{example}

\section{Infinite Geometric Series}

Finite geometric sums lead naturally to infinite series. Consider
$1+\tfrac12+\tfrac14+\tfrac18+\cdots$; its partial sums are $1$,
$\tfrac32$, $\tfrac74$, $\tfrac{15}{8}$, $\dots$, and these values appear
to approach $2$.

\begin{definition}
A sequence \emph{converges} if its terms approach a fixed limit. An
infinite series converges if its sequence of partial sums converges.
\end{definition}

\begin{theorem}
If $|r|<1$, then $\displaystyle\sum_{k=0}^\infty ar^k = \dfrac{a}{1-r}$.
\end{theorem}

\begin{proof}
The finite geometric formula gives $S_n = \dfrac{a(1-r^n)}{1-r}$. When
$|r|<1$, $r^n \to 0$ as $n\to\infty$, so $\lim_{n\to\infty} S_n =
\dfrac{a(1-0)}{1-r} = \dfrac{a}{1-r}$.
\end{proof}

\begin{example}
Evaluate $1+\tfrac12+\tfrac14+\tfrac18+\cdots$. Here $a=1$ and $r=1/2$,
so $\sum_{k=0}^\infty (1/2)^k = 1/(1-1/2) = 2$.
\end{example}

\begin{example}
The series $1+2+4+8+\cdots$ does not converge, because its partial sums
grow without bound; here $r=2$ fails the requirement $|r|<1$.
\end{example}

\section{The Complex Geometric Series}

The same algebra works when the ratio is a complex number.

\begin{theorem}
For any complex number $z \neq 1$, $1+z+z^2+\cdots+z^{n-1} =
\dfrac{1-z^n}{1-z}$.
\end{theorem}

\begin{proof}
Let $S_n = 1+z+z^2+\cdots+z^{n-1}$. Multiplying by $z$ gives $zS_n =
z+z^2+\cdots+z^n$. Subtracting, $S_n - zS_n = 1-z^n$, so $(1-z)S_n =
1-z^n$ and $S_n = \dfrac{1-z^n}{1-z}$.
\end{proof}

This proof is identical to the real-valued case; only the interpretation
changes. When $z = e^{i\theta}$, the successive terms $1, e^{i\theta},
e^{i2\theta}, e^{i3\theta}, \dots$ are unit vectors rotated by equal
angles, so the partial sums trace a polygonal path in the complex plane
rather than moving along the real line.

\begin{example}
Let $z = e^{2\pi i/3}$. Then $1+z+z^2 = \dfrac{1-z^3}{1-z} = 0$, exactly
the roots-of-unity result encountered in Chapter~\ref{ch:ch16-complex-euler}.
\end{example}

\begin{algorithmproc}{How to Identify and Sum an Arithmetic or Geometric Series}
Determine whether consecutive terms differ by a constant amount or are
multiplied by a constant ratio. Apply the arithmetic formulas when a
common difference exists, and the geometric formulas when a common ratio
exists. For an infinite geometric series, verify that $|r|<1$ before
applying $a/(1-r)$; if $|r|\ge1$, the series diverges and has no finite
sum.
\end{algorithmproc}

\begin{practicenow}
Compute $3+6+12+24+48$. Then determine whether $3+6+12+24+\cdots$
converges.
\end{practicenow}

\section{Connections to Earlier Chapters}

Several ideas encountered earlier in this book are instances of
geometric sequences. Repeated amplitwists satisfy $A^n = s^nR(n\theta)$
from Chapter~15: the scaling factor $s^n$ forms a geometric sequence,
while the angles $n\theta$ form an arithmetic sequence. Similarly,
$(re^{i\theta})^n = r^ne^{in\theta}$ from Chapter~16 contains the same
pair of patterns, with the magnitude evolving geometrically while the
angle evolves arithmetically. The concepts introduced in this chapter
therefore organize ideas that had already appeared, unnamed, in
transformation geometry and complex numbers.

\section{Applications}

Geometric sequences and series appear throughout mathematics and
science, in every setting where a quantity is repeatedly multiplied by a
fixed factor: compound interest, population growth and decay,
radioactive decay, recursive algorithms, and self-similar or fractal
structures all generate a geometric sequence at their core.

\begin{example}
An account earning $5\%$ annual interest, compounded once per year, has
balance $B_n = P(1.05)^n$ after $n$ years, where $P$ is the initial
deposit: a geometric sequence with common ratio $r = 1.05$. If \$1000 is
deposited and left untouched, the balance after $10$ years is $B_{10} =
1000(1.05)^{10} \approx 1628.89$. The total interest earned across those
ten years, summed year by year, is itself a partial sum of a related
geometric series, since each year's interest payment is $5\%$ of a
balance that has itself grown geometrically.
\end{example}

\section{Looking Ahead}

A finite geometric sum of complex exponentials produces a finite
superposition of rotating vectors. The central idea of Fourier analysis
is to combine many such rotations, each with its own frequency and
amplitude, to represent periodic phenomena.
Chapter~\ref{ch:ch19-fourier-intuition} asks what happens when
infinitely many such rotating components, rather than the single ratio
$z$ studied here, are combined to model waves, signals, and functions.

\section{Exercises}

\subsection*{Arithmetic Sequences and Series}
\begin{exercise}
Find the tenth term of the arithmetic sequence $4, 7, 10, 13, \dots$.
\end{exercise}

\begin{exercise}
Compute $5+9+13+\cdots+101$.
\end{exercise}

\begin{exercise}
Derive the arithmetic series formula independently, without referring
back to the proof given in the text.
\end{exercise}

\begin{exercise}
A theater has $20$ seats in its first row, with each subsequent row
having $2$ more seats than the row before it. Find the total number of
seats in the first $15$ rows.
\end{exercise}

\subsection*{Geometric Sequences and Series}
\begin{exercise}
Find the eighth term of $3, 6, 12, 24, \dots$.
\end{exercise}

\begin{exercise}
Evaluate $1+3+9+27+81$.
\end{exercise}

\begin{exercise}
Evaluate $\displaystyle\sum_{k=0}^\infty \left(\frac13\right)^k$.
\end{exercise}

\begin{exercise}
Determine whether $1-\tfrac12+\tfrac14-\tfrac18+\cdots$ converges, and
find its sum if it does.
\end{exercise}

\begin{exercise}
A ball is dropped from a height of $10$ meters, and after each bounce it
rises to $60\%$ of its previous height. Find the total vertical distance
the ball travels, accounting for infinitely many bounces.
\end{exercise}

\subsection*{Complex Geometric Series}
\begin{exercise}
Compute $1+i+i^2+i^3$.
\end{exercise}

\begin{exercise}
Verify directly that $1+e^{2\pi i/3}+e^{4\pi i/3}=0$.
\end{exercise}

\begin{exercise}
Use the geometric-series formula to derive the same result.
\end{exercise}

\begin{exercise}
Compute $1+z+z^2+z^3$ for $z=e^{i\pi/2}$, and interpret the partial sums
as vectors in the complex plane.
\end{exercise}

\subsection*{Mixed Problems}
\begin{exercise}
Determine whether the sequence $2, 5, 10, 17, 26, \dots$ is arithmetic,
geometric, or neither, and justify your answer.
\end{exercise}

\begin{exercise}
An initial deposit of \$2000 earns $4\%$ annual interest, compounded
once per year. Find the balance after $8$ years.
\end{exercise}

\begin{exercise}
Explain the difference between a sequence and a series, and give an
example of each.
\end{exercise}

\subsection*{Challenge Problems}
\begin{exercise}
Show that the scaling factors in repeated amplitwists form a geometric
sequence.
\end{exercise}

\begin{exercise}
Show that the angles in repeated amplitwists form an arithmetic
sequence.
\end{exercise}

\begin{exercise}
Explain why $(re^{i\theta})^n = r^ne^{in\theta}$ contains both a
geometric and an arithmetic progression simultaneously.
\end{exercise}

\begin{exercise}
Explain why a finite sum of complex exponentials may be interpreted as a
sum of rotating vectors.
\end{exercise}

\begin{exercise}
Prove that an infinite geometric series with $|r|\ge1$ cannot converge,
by examining what happens to the partial sums $S_n$ as $n\to\infty$.
\end{exercise}
